% Part: computability
% Chapter: computability-theory
% Section: non-comp-set

\documentclass[../../../include/open-logic-section]{subfiles}

\begin{document}

\olfileid{cmp}{thy}{ncp}
\olsection{存在不可计算的集合}

前面已经看到，每个可计算的集合都是可计算枚举的。其逆命题成立吗？下面将说明，一般而言，它并不成立。

\begin{thm}\ollabel{thm:K-0}
令 $K_0$ 为集合 $\Setabs{\tuple{e, x}}{\cfind{e}(x) \fdefined}$。
则 $K_0$ 是可计算枚举的，但不是可计算的。
\end{thm}

\begin{proof}
要看出 $K_0$ 是可计算枚举的，注意它是如下定义的函数~$f$ 的定义域：
\[
f(z) = \umin{y}{(\len{z} = 2 \land T((z)_0, (z)_1, y))}.
\]
因为，如果 $\cfind{e}(x)$ 有定义，$f(\tuple{e, x})$ 会找到一个停机计算序列；如果 $\cfind{e}(x)$ 无定义，那么 $f(\tuple{e, x})$ 也无定义；而如果 $z$ 甚至不编码一个数对，那么 $f(z)$ 同样无定义。

$K_0$ 不可计算这一事实正是停机问题的不可判定性，即 \olref[hlt]{thm:halting-problem}。
\end{proof}

集合 $K_0$ 是使得 $\cfind{e}(x) \fdefined$ 的数对 $\tuple{e,x}$ 的集合，即 $\tuple{e,x} \in K_0$ 当且仅当 $\cfind{e}$ 在输入~$x$ 上有定义（停机），因此它也被称为“停机集”。集合 $K = \Setabs{e}{\cfind{e}(e) \fdefined}$ 是“自停机集”。它常被用作标准的不可判定集。

\begin{thm}\ollabel{thm:K}
自停机集 $K = \Setabs{e}{\cfind{e}(e) \fdefined}$ 是可计算枚举的，但不是可判定的。
\end{thm}

\begin{proof}
假设 $K$ 是可判定的，即其特征函数 $\Char{K}$ 是可计算的。令
\[d(e) = \begin{cases}
  1 & \text{若\/ $\Char{K}(e) = 0$}\\
  \fundefined & \text{否则。}
  \end{cases}
  \]
设 $k$ 为~$d$ 的指标，即 $d \simeq \cfind{k}$。那么 $d(k)
  \simeq \cfind{k}(k)$。这与如下事实相矛盾：由 $d$ 的定义可知，$d(k) \fdefined$ 当且仅当 $\cfind{k}(k) \fundefined$。

$K$ 是 $f(x) = \umin{y}{T(x,x,y)}$ 的定义域，因此是可计算枚举的。
\end{proof}

\end{document}